\documentclass[11pt,letterpaper]{article}
\usepackage[T1]{fontenc}
\usepackage{lmodern,amsmath,amssymb,amsthm,microtype}
\usepackage[margin=1in]{geometry}
\usepackage[hidelinks]{hyperref}
\providecommand{\tightlist}{\setlength{\itemsep}{0pt}\setlength{\parskip}{0pt}}
\DeclareUnicodeCharacter{2264}{\ensuremath{\le}}
\DeclareUnicodeCharacter{2265}{\ensuremath{\ge}}
\DeclareUnicodeCharacter{2260}{\ensuremath{\ne}}
\DeclareUnicodeCharacter{2212}{\ensuremath{-}}

\DeclareUnicodeCharacter{2211}{\ensuremath{\sum}}

\title{Six grid additions to \(G_{449}\): the four-rim one-neighbor case}
\author{Aafi Khan Malek\\ \small Independent researcher, Ahmedabad, India\\
\small\href{mailto:aafimalek2032@gmail.com}{aafimalek2032@gmail.com}}
\date{28 November 2026}
\begin{document}
\maketitle
\textbf{Theorem.} Let six distinct points be added to the
Heule--Przybocki graph \(G_{449}\) in its denominator-seven grid.
If exactly one new point has one original neighbor and it is adjacent
to exactly four of the other five new points, the extension remains
five-colorable. We use the earlier verified five-point extension
theorem and a finite exact enumeration. This is not a new bound for
the entire plane and is not a Lean theorem.

Assume every five-point extension of the pinned G449 graph is
five-colorable. For an added point \(v\), let \(b(v)\) be its number of
original neighbors and \(d_H(v)\) its degree among six additions. A
failing configuration has \[
\delta_{\rm add}=\min_v\left(b(v)+d_H(v)\right)\ge5,
\] by deletion and greedy extension. Here the sole \(b=1\) point \(h\)
has exactly four new neighbors \(R\), and misses a fifth point \(w\). A
\(b=0\) point would have to be universal and share at least three
neighbors with \(h\), so every point other than \(h\) has \(b\ge2\).

\section{Complete finite search}\label{complete-finite-search}

\begin{enumerate}
\def\labelenumi{\arabic{enumi}.}
\tightlist
\item
  Regenerate the 7,056 grid points with \(b\ge2\) and 222,264 with
  \(b=1\) from the exact signed step catalogue and pinned base
  coordinates. For every possible rim point translate by all 546 steps.
  Partition the 3,852,576 generated incidences by the exact hub key and
  keep hubs with \(b=1\), storing the index of that original neighbor.
  The prior independently checked incidence histogram gives 109,683
  retained incidences, 30,233 hubs with any available rim points, 12,558
  with at least four, and 872,165 raw four-rim choices.
\item
  For each hub, the rim sits on the unit circle and exact rim adjacency
  follows \(\pm60^\circ\) orbits, so components have at most six
  vertices. A selected rim point requires \(b(v)+d_R(v)\ge3\), because
  its hub edge and at most one edge to \(w\) contribute at most two.
  Iteratively peel vertices that cannot meet this even with all
  remaining rim neighbors. Combine eligible subsets from distinct
  components in fixed component order to obtain each four-rim set once.
\item
  Define the required rim contacts to be those with \(b(v)+d_R(v)=3\).
  The missed point must hit all such contacts. Because \(h,w\) are
  distinct and share at most two unit neighbors, \(t=|N(w)\cap R|\) is
  0, 1, or 2. More than two required contacts renders this rim set
  impossible. A completion must satisfy \(b(w)\ge5-t\). Check the three
  cases separately:

  \begin{itemize}
  \tightlist
  \item
    \textbf{Zero.} No required contacts, \(b(w)\ge5\), and \(w\) is not
    adjacent to any rim or to \(h\), or equal to any selected point.
    Static bitsets range over the 574 \(b\ge5\) points. Subtract the hub
    neighbors and the rim points and their neighbors; this counts every
    valid isolated completion exactly once. For a given proper base
    palette, the bitset of isolated sites with a nonempty list is
    precomputed. The union of those bitsets over palettes that color the
    hub plus its four rims gives exactly the covered isolated
    completions.
  \item
    \textbf{One or two.} Enumerate the union of neighbor sets of the
    four rims restricted to \(b\ge3\). For a candidate \(w\), compute
    the complete four-bit contact mask by exact step lookup; visit it
    only from the least contacted rim index. Reject if it sees \(h\),
    equals a selected point, misses a required contact, or has
    \(b(w)+t<5\). This produces each valid one-contact or two-contact
    completion exactly once.
  \end{itemize}
\item
  For every coloring palette, remove the color of the hub's sole
  original neighbor from the hub list. A trial hub color is then
  forbidden at all four rim vertices but \textbf{not} at \(w\). The
  exact finite list search checks all rim edges, all edges to \(w\), and
  all lists derived from the original neighbors. An optional shortcut
  rejects a palette if it cannot already color the hub and its four
  rims. Residuals are bank gaps, not witnesses of six-chromaticity.
\end{enumerate}

Uniqueness of the hub, of the four-rim set, and of the missed vertex
makes this duplicate-free globally. The only other multiplicity in
generation (multiple contacts in the attached branch) is removed by the
least-contact rule. Each whole hub is in a single shard.

\section{Checked totals}
We replayed all 32 disjoint shards under the same memory guard.
The catalogue generated 872,165 raw four-rim sets, of which 83,818
passed the relaxed degree condition. The exact eligible completions
are 3,833,277 with no missed-rim contact, 145,271 with one, and
33,796 with two. All 4,012,344 are five-colorable using the verified
palette bank. The residual set is empty. A fresh independent run
matched all shard receipts exactly.
\par
Together with the previous low-degree and off-grid verifications,
this leaves only all-grid six-point additions in which each new
point has at least two original neighbors. Those cases are not
addressed by this proof.

\end{document}
